\section{Adaptationism: local formal dependencies}
\label{sec:p11-rat106-proof-annex}

\begin{annexitem}{annex:p11.rat106-lottery}{thm:p11-rat106-lottery}{Disease lotteries and the threshold comparison}
\begin{proof}
The expectation of the certain number saved is 200. The lottery's
expectation is \(600/3+2\cdot0/3=200\). Expected utility instead
compares \(u(200)\) with \(u(600)/3+2u(0)/3\), so equality of
these quantities is both necessary and sufficient for indifference
under the stipulated expected-utility criterion. If
\(u(x)=ax+t\), both values are \(200a+t\), proving the affine
case. Equal means do not force equality for arbitrary utility:
on nonnegative outcomes \(u(x)=x^2\) gives 40,000 for the sure
option and 120,000 for the lottery, while \(u(x)=\sqrt{x}\)
gives \(\sqrt{200}>\sqrt{600}/3\) for the sure option. The last
strict inequality follows by squaring nonnegative terms: \(200>600/9\).
These examples are increasing utilities on the relevant domain and
give strict opposite rankings.

At threshold zero every possible outcome succeeds, so both success
probabilities are one. For \(0<k\leq200\), the sure option
succeeds with probability one and the lottery only on its 600 branch,
with probability one third. For \(200<k\leq600\), the sure
option fails and the lottery still succeeds on the 600 branch.
Those intervals include the boundary values 200 and 600 exactly as
stated and exhaust the permitted domain. Maximizing these probabilities
therefore gives the asserted comparisons, and conversely their strict
rankings force selection of the stated option under that criterion.

The deaths description replaces each saved number \(x\) by
\(600-x\). It is a bijection on the outcome descriptions and
does not change any branch probability. Transporting the same
utility by \(\widetilde u(y)=u(600-y)\), or transporting the
threshold event to \(y\leq600-k\), leaves every expected
utility or success probability unchanged. A changed ranking then
requires a changed interpretation or criterion, not the relabeling
itself. No assertion about participants' actual interpretations is
used in this proof.
\end{proof}
\end{annexitem}

\begin{annexitem}{annex:p11.rat106-conditional}{thm:p11-rat106-conditional}{Falsifiers and modus tollens}
\begin{proof}
Classical material implication \(P\to Q\) is evaluated as
\(\neg P\lor Q\). It is false if and only if both disjuncts
are false, which is equivalent to \(P\) true and \(Q\)
false. This proves both directions of the falsifier characterization.
If the conditional is true while \(Q\) is false, its second
disjunct is false and its first must be true; hence \(\neg P\).
Equivalently a true \(P\) would create the just-characterized
falsifier and contradict the true conditional.

For information search, a case known to have \(P\) can have
an undisclosed false \(Q\), and a case known to have
\(\neg Q\) can have an undisclosed true \(P\); each
therefore can conceal the sole falsifying combination. A case
known to have \(\neg P\) makes the first disjunct true
regardless of the undisclosed side. A case known to have
\(Q\) makes the second disjunct true. Neither can conceal a
falsifier. Thus inspecting all and only the first two types suffices
to find every possible falsifier in the specified four-type task.
Omitting either type fails that guarantee because an otherwise
uninspected case of its type can have the falsifying undisclosed
side. The guarantee is about logical possibility and the given
one-known-side setup, not utility, a sampling frequency, the
probability that a counterexample exists, or deontic culpability.
\end{proof}
\end{annexitem}

\begin{annexitem}{annex:p11.rat106-best-response}{thm:p11-rat106-best-response}{The supplement's S1--S2 and the best response}
\begin{proof}
The law of total probability gives \(s_i\) and \(r_i\)
as the conditional cooperation probabilities given each actual dyad
type. Conditional independence makes the mutual-cooperation
probability \(s_1s_2\) in one-shot dyads and \(r_1r_2\)
in repeated dyads. In a one-shot round player one's payoff is
\(b\) times its partner's cooperative action minus \(c\)
times its own, so its conditional expectation is \(bs_2-cs_1\).

In a repeated dyad, mutual cooperation on every round has expected
length \(1/(1-w)\). Indeed the probability of reaching round
\(n\geq1\) is \(w^{n-1}\), and the nonnegative series
\(\sum_{n\geq1}w^{n-1}\) converges to that value because
\(0\leq w<1\). Expected mutual-cooperation payoff is
\((b-c)/(1-w)\). If only player one cooperates initially,
the first-round payoff is \(-c\); TIT-for-TAT defects after
the partner's initial defection and subsequent payoffs are zero.
The opposite pairing gives \(b\), and mutual defection gives
zero. Consequently the repeated conditional payoff is
\[
 r_1r_2\frac{b-c}{1-w}-r_1(1-r_2)c+(1-r_1)r_2b
 =br_2-cr_1+\frac{w(b-c)}{1-w}r_1r_2.
\]
Weighting the two conditional payoffs by \(P\) and \(1-P\)
proves the stated \(V_1\). Substitution of the definitions
of \(s_i,r_i\) expands it into the four belief-pair branches
for each actual dyad type, with the same branch probabilities
and payoffs as the supplement's S1. Thus this is an algebraically
equivalent expression, not a different game.

Only \(s_1\) and \(r_1\) vary with \(\alpha_1\):
their respective coefficients are \(\Theta_{\mathrm{OS}}\)
and \(\Theta_R\). Holding every other parameter fixed gives
the displayed coefficient by direct differentiation, or by
collecting the linear terms. For \(0<w<1\),
\(w/(1-w)=1/((1/w)-1)\), recovering the supplement's S2.
The first expression, unlike the latter, also defines the value
zero at \(w=0\).

An affine function on \([0,1]\) with positive slope is strictly
increasing, and so has its unique maximum at one; with negative
slope it has its unique maximum at zero; with zero slope it is
constant. These conditions are also necessary: the value at the
opposite endpoint or an interior value disproves a claimed unique
maximum under either other sign. This proves the complete
best-response classification, including all ties.

If \(w,(1-P),\Theta_R,r_2\) are positive, their product
divided by \(1-w\) is positive. The positive term is that
fixed coefficient times \(b-c\), whereas the subtracted cost
is fixed and finite. Taking \(b-c\) sufficiently large
therefore makes the slope positive. If \(w=0\),
\(\Theta_R=0\), or \(r_2=0\), that positive term
vanishes and the slope is nonpositive; it is strictly negative
when \(P\Theta_{\mathrm{OS}}+(1-P)\Theta_R>0\).
These are counterexamples to dispensing with the positive-factor
hypotheses. Neither the linear maximization nor the preceding
payoff calculation proves convergence under mutation and selection.
\end{proof}
\end{annexitem}
